\documentclass{article}

\usepackage{amsmath,amssymb,amsthm,mathtools}
\usepackage[margin=0.95in]{geometry}
\usepackage{enumitem}
\usepackage{microtype}
\usepackage[dvipsnames]{xcolor}
\usepackage[pdftex,colorlinks,backref=page,citecolor=blue]{hyperref}

\allowdisplaybreaks
\setlist[enumerate]{itemsep=2pt,topsep=4pt,leftmargin=*}

\theoremstyle{plain}
\newtheorem{lemma}{Lemma}
\newtheorem{theorem}{Theorem}
\newtheorem{proposition}{Proposition}
\newtheorem{example}{Example}
\newtheorem{corollary}{Corollary}
\newtheorem{conjecture}{Conjecture}
\theoremstyle{definition}
\newtheorem{remark}{Remark}
\newtheorem{definition}{Definition}
\newcommand{\E}{\mathbb{E}}

% ============= COLOR SYSTEM =============
\definecolor{provedgreen}{RGB}{20,120,40}
\definecolor{openred}{RGB}{180,30,30}
\newcommand{\proved}{\textcolor{provedgreen}{\textbf{[Proved]}}\,}
\newcommand{\condcolor}{\textcolor{provedgreen}{\textbf{[Proved given Hypothesis~H]}}\,}
\newcommand{\open}{\textcolor{openred}{\textbf{[Open / numerical]}}\,}
% =========================================

\title{Heterogeneity and Incentives in Dynamic Matching Markets:\\
Outside Options, Liquidity, and Competition between Exchanges}
\author{Sinaya Joshi \\ \textit{Stanford University}}

\begin{document}
\maketitle

\begin{abstract}
We extend the two-type continuous-time dynamic matching model to allow
easy-to-match (E) agents to face outside-option pressure: E agents become
critical (depart) at rate $\rho \ge 1$ rather than rate 1.  This models
altruistic donors or versatile workers who can exit the market when attractive
alternatives appear elsewhere.  We derive closed-form loss functions for all
four canonical policies (Greedy, Patient, Type-Aware Greedy, Type-Aware
Patient) as functions of $\rho$ and $d_1$ (H-E compatibility density),
revealing a fundamental asymmetry: \textbf{Patient-type policies are directly
exposed to outside-option pressure while Greedy-type policies are not, at
leading order.}

Key results: $L^{\text{Patient}}(d_1,\rho)=4\rho W(d_1/(8\rho))/d_1$ and
$L^{\text{TA-Greedy}}(d_1,\rho)=2W(d_1/4)/d_1$ (independent of $\rho$),
yielding the duality $L^{\text{Patient}}(d_1,\rho)=L^{\text{TA-Greedy}}(d_1/(2\rho))$.
For TA-Patient the exponential decay rate becomes $e^{-d_1/(4(1+\rho))}$:
stronger outside options compress the exponential advantage.  We analyze
competition between platforms, showing private exchanges can dominate a
public monopoly when markets are thick enough.  These findings speak to the
empirical debate over kidney-exchange consolidation.
\end{abstract}

\tableofcontents

%----------------------------------------------------------------------
\section{Introduction}
\label{sec:intro}

A central question in matching-market design is whether to centralize or
fragment.  In kidney exchange, the US relies on multiple hospital-run
consortia---the National Kidney Registry (NKR, daily clearings), the
Alliance for Paired Kidney Donation (daily or on new-entry), and the OPTN
clearinghouse (weekly)---alongside hospital-specific programs that match
locally before accessing national pools.  The economic literature has
emphasized gains from thicker markets: pooling more agents raises match
rates, especially for hard-to-match types.  Yet fragmentation persists and
in some cases improves outcomes.

This paper identifies a mechanism for why: when easy-to-match agents---altruistic
donors, bridge donors, versatile workers---face outside options that cause
them to leave the market, the key design variable is not pool size alone but
how quickly the market captures value from E agents before they exit.  A
platform that credibly reduces easy agents' effective departure rate (through
priority credits, better match quality, or commitment mechanisms) delivers
strictly higher welfare for hard-to-match agents, and private platforms
unconstrained by public procurement rules may have a comparative advantage
on this dimension.

\paragraph{The heterogeneous model.}
We model the market as a continuous-time stochastic process with two agent
types: hard-to-match (H) and easy-to-match (E).  The compatibility graph on
the pool is a Stochastic Block Model: H-H edges are forbidden, H-E edges
appear independently with probability $d_1/(2m)$, and E-E edges with
probability $d_2/(2m)$, where $d_2>d_1>0$ captures heterogeneity.  This is
the \emph{sparse} regime: compatibility probabilities vanish as $1/m$, so
the expected H-E degree of a pooled agent is $O(1)$, unlike the dense regime
of Ashlagi, Nikzad, and Strack (2023) where $p_{HE}$ is a fixed constant.
Empirically, Liu et al.\ (2018) find in DiDi carpooling data that moderate
heterogeneity favors Patient, but pronounced heterogeneity can reverse this
advantage---consistent with the crossover behavior our model predicts.

\paragraph{The key parameter: $\rho$.}
We introduce the E departure rate $\rho\ge 1$.  When $\rho=1$, E agents
depart no faster than H agents (baseline).  When $\rho>1$, E agents become
critical faster.  The steady-state E pool size is proportional to $1/\rho$,
so the inventory of potential rescue partners for H agents shrinks as $\rho$
rises.

\paragraph{Main findings.}
\begin{enumerate}
\item \emph{Greedy policies are $\rho$-insensitive at leading order.}
  Neither Greedy nor TA-Greedy loss depends on $\rho$ as $d_2\to\infty$.
  The H-E density $d_1$ is the binding constraint.
\item \emph{Patient loss has an exact closed form in $\rho$:}
  $L^{\text{Patient}}(d_1,\rho)=4\rho W(d_1/(8\rho))/d_1$, where $W$ is
  the Lambert $W$ function.
\item \emph{TA-Patient's exponential advantage erodes with $\rho$.}
  The decay exponent decreases from $d_1/8$ (at $\rho=1$) to $d_1/(4(1+\rho))$.
\item \emph{Duality.}
  $L^{\text{Patient}}(d_1,\rho)=L^{\text{TA-Greedy}}(d_1/(2\rho))$:
  outside-option pressure on E is equivalent to a rescaling of H-E density.
\item \emph{Competition can be welfare-improving.}
  A private platform that uses TA-Patient and reduces $\rho$ dominates a
  public Patient platform in thick markets.
\end{enumerate}

\begin{remark}[Notation conflict]
In prior work on asymmetric markets, $\rho$ sometimes denotes the arrival-rate
ratio between types.  Here $\rho$ is exclusively the E departure rate; arrival
rates remain equal at $m/2$ each.
\end{remark}

%----------------------------------------------------------------------
\section{Model}
\label{sec:model}

\subsection{Setup}

The market has two types: hard-to-match (H) and easy-to-match (E).  Both
arrive as Poisson processes at rate $m/2$ each (total rate $m$).

\paragraph{Compatibility.}
Any two agents in the pool are independently compatible with type-pair
probabilities:
\[
  p_{HH}=0,\quad
  p_{HE}=p_{EH}=\frac{d_1}{2m},\quad
  p_{EE}=\frac{d_2}{2m},\qquad d_2>d_1>0.
\]

\paragraph{Departure rates.}
H agents become critical at rate 1; E agents become critical at rate
$\rho\ge 1$.  A critical agent must be matched immediately or it perishes.

\paragraph{Pool state.}
Let $N_t=(H_t,E_t)$ denote the pool state.  In rescaled variables
$x=H_t/m$ and $y=E_t/m$, the pool evolves via arrivals, criticalities, and
matches.

\paragraph{Neighbor probabilities.}
Given pool state $(h,e)$, the probability that a given H agent has at least
one compatible (E) neighbor is:
\[
  \pi_H(h,e)
    =1-\Bigl(1-\tfrac{d_1}{2m}\Bigr)^{\!e}
    \approx 1-e^{-d_1 y/2},
\]
and the probability that a given E agent has at least one compatible (H or E)
neighbor is:
\[
  \pi_E(h,e)
    =1-\Bigl(1-\tfrac{d_1}{2m}\Bigr)^{\!h}\Bigl(1-\tfrac{d_2}{2m}\Bigr)^{\!e}
    \approx 1-e^{-(d_1 x+d_2 y)/2}.
\]

\subsection{Four policies}
\label{sec:policies}

We study four canonical policies:
\begin{itemize}
\item \textbf{Greedy}: match on arrival, choosing uniformly among compatible
  neighbors in the pool.
\item \textbf{TA-Greedy}: match on arrival, prioritizing H-type neighbors for
  E arrivals.
\item \textbf{Patient}: match at criticality, choosing uniformly among
  compatible neighbors in the pool.
\item \textbf{TA-Patient}: match at criticality, prioritizing H-type neighbors
  for E criticals.
\end{itemize}
Under any Greedy variant the pool is edgeless: arriving agents have not yet
encountered each other, so no two pooled agents have a tested edge.  Under
any Patient variant the pool retains a random graph on pooled agents.

A fifth policy, \textbf{TA-Hybrid} (match H immediately upon receiving a
compatible E neighbor; be patient otherwise), interpolates between TA-Greedy
and TA-Patient and is left for future analysis.

\subsection{Drift and mean-field fixed point}

For a given policy, the drift of $(H_t,E_t)$ is:
\[
  f_H(h,e)=\E[\Delta H_t\mid H_t=h,E_t=e],\qquad
  f_E(h,e)=\E[\Delta E_t\mid H_t=h,E_t=e].
\]
A \emph{mean-field fixed point} $(x^*,y^*)$ satisfies
$f_H(mx^*,my^*)=0$ and $f_E(mx^*,my^*)=0$.  Under standard concentration
arguments (supported numerically by occupancy heatmaps that show the
stationary distribution concentrating tightly near $(h^*,e^*)$), the
long-run loss is well approximated by evaluating at the fixed point.

\subsection{Loss}

The long-run loss is the fraction of agents who perish unmatched:
\begin{equation}\label{eq:loss-def}
  L
  =\lim_{T\to\infty}
    \frac{\E[\text{agents perishing unmatched by }T]}
         {\E[\text{total arrivals by }T]}.
\end{equation}
In mean-field terms, at the fixed point $(x^*,y^*)$:
\begin{equation}\label{eq:loss-mf}
  L = x^*(1-\pi_H^*) + \rho\,y^*(1-\pi_E^*),
\end{equation}
where $\pi_H^*=1-e^{-d_1 y^*/2}$ and $\pi_E^*=1-e^{-(d_1 x^*+d_2 y^*)/2}$.
The two terms are the H-perish rate (H becomes critical without a compatible
neighbor) and the E-perish rate (E becomes critical at rate $\rho$ without
any compatible neighbor).

%----------------------------------------------------------------------
\section{Analysis}
\label{sec:mf}

\subsection{Greedy: leading-order independence from $\rho$}
\label{sec:greedy-rho}

Under Greedy, the drift of the rescaled pool $(x,y)$ is:
\begin{align}
  F_H(x,y) &= \tfrac{1}{2}(1-\pi_H) - x - \tfrac{1}{2}\,r_{EH},
    \label{eq:greedy-drift-H}\\
  F_E(x,y) &= \tfrac{1}{2}(1-\pi_E) - \rho y - \tfrac{1}{2}\,\pi_H
               - \tfrac{1}{2}\,r_{EE},
    \label{eq:greedy-drift-E}
\end{align}
where $r_{EH}$ and $r_{EE}$ are the probabilities that an arriving E agent
matches an H or E agent, respectively.

\begin{lemma}[Greedy loss is $\rho$-independent at leading order]
\label{lem:greedy-rho-indep}
\proved As $d_2\to\infty$, the Greedy fixed-point $x^*$ is the unique root
in $(0,1/2)$ of
\[
  g_{d_1}(x)
    :=x\exp\!\left(\frac{d_1 x(1-x)}{2(1-2x)}\right)=1,
\]
independent of $\rho$.  Hence $L^{\mathrm{Greedy}}(d_1,\rho)=x^*(d_1)+O(1/d_2)$.
\end{lemma}

\begin{proof}
In rescaled-degree variables with $\pi_H=1-e^{-d_1 y/2}$ and similarly for
$\pi_E$ and $r_{EE}$, the fixed-point equations become:
\begin{align}
  x+\rho y
    &= e^{-d_1 y/2}+e^{-(d_1 x+d_2 y)/2}-1, \tag{S1}\\
  x-\rho y
    &= \bigl(1-e^{-(d_1 x+d_2 y)/2}\bigr)\cdot\frac{d_2 y}{d_1 x+d_2 y}.
    \tag{S2}
\end{align}

\emph{Step 1: Scaling of $y$.}
Write $\rho y=c/d_2$ for $c=c(x,d_1)>0$.  Substituting into (S2) and
dropping the $O(1/d_2)$ term on the left side:
\[
  x = \bigl(1-e^{-(d_1 x+c)/2}\bigr)\cdot\frac{c}{d_1 x+c}.
\]
Setting $\alpha=d_1 x+c$ and rearranging: $x\alpha=c(1-e^{-\alpha/2})$,
which expands to $d_1 x^2+xc=c-ce^{-\alpha/2}$.  At leading order (the
exponential term is $O(e^{-c/2})$ and $c\to\infty$ as $d_2\to\infty$):
\begin{equation}\label{eq:c-from-x}
  c = \frac{d_1 x^2}{1-2x}. \tag{$*$}
\end{equation}
This relation is independent of $\rho$.

\emph{Step 2: Scalar equation for $x$.}
Substituting $y=c/(d_2\rho)$ into (S1) and taking $d_2\to\infty$
(so $e^{-d_1 y/2}\to 1$ and $e^{-(d_1 x+c)/2}$ is the surviving
exponential):
\[
  x \;\longrightarrow\; e^{-(d_1 x+c)/2}.
\]
Substituting \eqref{eq:c-from-x}:
\[
  x = \exp\!\left(-\frac{d_1 x}{2}-\frac{d_1 x^2}{2(1-2x)}\right)
     = \exp\!\left(-\frac{d_1 x(1-x)}{2(1-2x)}\right),
\]
which is $g_{d_1}(x)=1$, independent of $\rho$.

\emph{Step 3: Loss.}
Since $\rho y=c/d_2\to 0$ and $1-\pi_E=e^{-(d_1 x+d_2 y)/2}\to 0$, we
have $\rho y(1-\pi_E)=O(1/d_2)$.  Thus
$L=x(1-\pi_H)+O(1/d_2)=x^*(d_1)+O(1/d_2)$, independent of $\rho$.
\end{proof}

\begin{remark}[Intuition]
Under Greedy, matching happens at agent arrival.  When $d_2$ is large, the
E pool is dense enough to serve as a reliable standing inventory for H
arrivals regardless of how quickly E agents depart.  The binding constraint
is the H-E compatibility density $d_1$; the E departure rate is a
second-order correction.
\end{remark}

\subsection{TA-Greedy: same leading-order independence}
\label{sec:TAG-rho}

Under TA-Greedy, an arriving E agent with both H and E compatible neighbors
prioritizes the H neighbor.  At leading order as $d_2\to\infty$, the
fixed-point analysis runs as in Lemma~\ref{lem:greedy-rho-indep} with one
modification: the H-E compatibility density is effectively doubled (each E
arrival rescues the H neighbor immediately), which shifts the equation for
$x^*$.  The outcome is:
\[
  L^{\text{TA-Greedy}}(d_1,\rho)
    = \frac{2W(d_1/4)}{d_1}+O(1/d_2),
\]
independent of $\rho$ at leading order.

\emph{Sketch.}  The key step is identical to Step~1 in the Greedy proof: the
E pool satisfies $\rho y=O(1/d_2)$, so $\rho$ enters the fixed-point
equations only through the vanishing term $\rho y=c/d_2$.  The scalar
equation for $x^*$ that survives the $d_2\to\infty$ limit is $\rho$-free.
The Lambert-$W$ formula $2W(d_1/4)/d_1$ comes from applying the same
argument with H-priority, which changes the exponent in $g_{d_1}$ by a
factor of 2.

\subsection{Patient: closed-form $\rho$-dependence}
\label{sec:patient-rho}

Under Patient, agents wait in the pool until critical, then match uniformly
among compatible neighbors.

\begin{lemma}[Patient loss with $\rho$]
\label{lem:patient-rho}
\proved As $d_2\to\infty$, the Patient fixed point satisfies $x^*\to 1/2$
and $y^*=y^*(d_1,\rho)$ where $e^{-d_1 y^*/2}=4\rho y^*$.  The loss is:
\[
  \boxed{L^{\mathrm{Patient}}(d_1,\rho)
    = \frac{4\rho\,W(d_1/(8\rho))}{d_1}.}
\]
\end{lemma}

\begin{proof}
At the mean-field fixed point, the steady-state conditions for the H and E
pool sizes are:
\begin{align}
  \tfrac{1}{2}
    &= x + \rho y\,\pi_E\,\rho_{E\to H},
    \tag{P1}\\
  \tfrac{1}{2}
    &= x\,\pi_H + \rho y\bigl(1+\pi_E\,\rho_{E\to E}\bigr).
    \tag{P2}
\end{align}
In (P1): H flows in at rate $1/2$; it leaves at rate $x$ (all critical H,
whether matched or not) plus rate $\rho y\,\pi_E\rho_{E\to H}$ (H matched by
a critical E).  In (P2): E flows in at rate $1/2$; it leaves at rate
$x\pi_H$ (E matched by critical H) plus rate $\rho y(1+\pi_E\rho_{E\to E})$
(each critical E removes itself, and with probability $\pi_E\rho_{E\to E}$
it also removes a second E as the matched partner).

\emph{Step 1: $x\to 1/2$ as $d_2\to\infty$.}
As $d_2\to\infty$ with a positive E pool, critical E agents have almost
surely at least one compatible neighbor, but the fraction of those neighbors
that are H-type vanishes: $d_1 x/(d_1 x+d_2 y)\to 0$.  Under type-unaware
Patient, E matches uniformly, so $\pi_E\rho_{E\to H}\to 0$.  From (P1):
$1/2=x+o(1)$, hence $x\to 1/2$.

\emph{Step 2: Scalar equation for $y$.}
With $d_2\to\infty$, every critical E almost surely finds an E-compatible
neighbor, so $\pi_E\rho_{E\to E}\to 1$ and thus $1+\pi_E\rho_{E\to E}\to 2$.
Also $\pi_H=1-e^{-d_1 y/2}$.  Substituting $x=1/2$ in (P2):
\[
  \tfrac{1}{2}
    = \tfrac{1}{2}\bigl(1-e^{-d_1 y/2}\bigr) + 2\rho y.
\]
Simplifying:
\begin{equation}\label{eq:patient-transcendental}
  e^{-d_1 y/2} = 4\rho y.
\end{equation}

\emph{Step 3: Lambert~$W$ form.}
Multiply both sides of \eqref{eq:patient-transcendental} by $e^{d_1 y/2}$:
$1=4\rho y\,e^{d_1 y/2}$.  Set $u=d_1 y/2$, so $y=2u/d_1$:
\[
  1 = \frac{8\rho}{d_1}\,u\,e^u
  \implies
  u\,e^u = \frac{d_1}{8\rho}.
\]
By definition of the Lambert $W$ function, $u=W(d_1/(8\rho))$, giving
$y^*=2W(d_1/(8\rho))/d_1$.

\emph{Step 4: Loss formula.}
Adding (P1) and (P2) and using $\rho_{E\to H}+\rho_{E\to E}=1$ (a critical
E with at least one compatible neighbor matches exactly one agent), so that
$\pi_E\rho_{E\to H}+\pi_E\rho_{E\to E}=\pi_E$:
\[
  1 = x(1+\pi_H)+\rho y(1+\pi_E).
\]
From $L=x(1-\pi_H)+\rho y(1-\pi_E)$ and the equation above:
\[
  L = (x+\rho y)-(x\pi_H+\rho y\pi_E)
    = (x+\rho y)-\bigl(1-(x+\rho y)\bigr) = 2(x+\rho y)-1.
\]
With $x=1/2$ and $y^*=2W(d_1/(8\rho))/d_1$:
\[
  L = 2\!\left(\tfrac{1}{2}+\rho\cdot\frac{2W(d_1/(8\rho))}{d_1}\right)-1
    = \frac{4\rho\,W(d_1/(8\rho))}{d_1}. \qedhere
\]
\end{proof}

\begin{remark}[Behavior in $\rho$]
As $\rho\to 1$: recover the baseline Patient loss.
As $\rho\to\infty$: the argument of $W$ tends to 0, so
$W(d_1/(8\rho))\approx d_1/(8\rho)$, giving $L\to 1/2$ (half of all agents
perish).
As $\rho\to 0$: E agents are permanently sticky and $L\to 0$.
The log-elasticity is
$\partial\log L/\partial\log\rho = W(d_1/(8\rho))/(1+W(d_1/(8\rho)))\in(0,1)$.
\end{remark}

\subsection{TA-Patient: identities and exponential decay}
\label{sec:TAP-rho}

Under TA-Patient, critical E agents match their H-compatible neighbor first;
only if no H neighbor exists do they match an E.

\begin{lemma}[TA-Patient fixed-point identities]
\label{lem:TAP-identities-rho}
\proved Any solution $(x,y)$ of the TA-Patient fixed-point system
(as $d_2\to\infty$) satisfies:
\begin{align}
  L^{\mathrm{TA\text{-}Patient}} &= x\,e^{-d_1 y/2},
    \label{eq:TAP-loss-identity}\\
  x\,e^{-d_1 y/2} &= 2\rho y\,e^{-d_1 x/2}.
    \label{eq:TAP-cross-identity}
\end{align}
\end{lemma}

\begin{proof}
As $d_2\to\infty$, the TA-Patient fixed-point equations are:
\begin{align}
  \tfrac{1}{2} &= x + \rho y\bigl(1-e^{-d_1 x/2}\bigr), \tag{TAP1}\\
  \tfrac{1}{2} &= x\bigl(1-e^{-d_1 y/2}\bigr)
                  +\rho y\bigl(1+e^{-d_1 x/2}\bigr). \tag{TAP2}
\end{align}
In (TAP1): all critical H leave at rate $x$; critical E's match H with
probability $1-e^{-d_1 x/2}$ (probability of having at least one H neighbor)
at rate $\rho y$.  In (TAP2): critical H's match E with probability
$\pi_H=1-e^{-d_1 y/2}$ at rate $x$; critical E's always leave at rate
$\rho y$ (the $+1$ term), and additionally pull a second E from the pool
when they find no H neighbor and instead match an E (the $+e^{-d_1 x/2}$
term, since $d_2\to\infty$ guarantees an E neighbor whenever needed).

\emph{Step 1: Summing (TAP1) and (TAP2).}
\begin{align*}
  1
    &= \bigl[x+\rho y(1-e^{-d_1 x/2})\bigr]
       +\bigl[x(1-e^{-d_1 y/2})+\rho y(1+e^{-d_1 x/2})\bigr]\\
    &= 2x+2\rho y - xe^{-d_1 y/2}
      \underbrace{-\rho ye^{-d_1 x/2}+\rho ye^{-d_1 x/2}}_{=\,0}.
\end{align*}
Hence $L=2(x+\rho y)-1=xe^{-d_1 y/2}$ (using the same aggregate-balance
identity as in the Patient proof, which holds because TA-Patient also
satisfies $\pi_E\rho_{E\to H}+\pi_E\rho_{E\to E}=\pi_E$).

\emph{Step 2: Subtracting (TAP1) from (TAP2).}
\begin{align*}
  0
    &= \bigl[x(1-e^{-d_1 y/2})+\rho y(1+e^{-d_1 x/2})\bigr]
       -\bigl[x+\rho y(1-e^{-d_1 x/2})\bigr]\\
    &= -xe^{-d_1 y/2}+2\rho ye^{-d_1 x/2}.
\end{align*}
Thus $xe^{-d_1 y/2}=2\rho ye^{-d_1 x/2}$.
\end{proof}

\begin{conjecture}[TA-Patient exponential decay with $\rho$]
\label{conj:TAP-rho}
For $d_1\to\infty$ at fixed $\rho\ge 1$:
\[
  \boxed{L^{\mathrm{TA\text{-}Patient}}(d_1,\rho)
    \;\sim\; \frac{(2\rho)^{1/(1+\rho)}}{2(1+\rho)}\,e^{-d_1/(4(1+\rho))}.}
\]
\end{conjecture}

The conjecture is supported by numerical simulation.  A heuristic derivation:
equations \eqref{eq:TAP-loss-identity}--\eqref{eq:TAP-cross-identity} together
with (TAP1) suggest the saddle-point scaling $y/x\to 1/\rho$ for large $d_1$,
leading to the linear system $d_1/(4(1+\rho))$ in the exponent and the
prefactor $(2\rho)^{1/(1+\rho)}/(2(1+\rho))$ from matching boundary
conditions.

The decay exponent $A(\rho)=1/(4(1+\rho))$ is strictly decreasing in $\rho$:
\begin{center}
\begin{tabular}{r|cccccc}
$\rho$   & $0.5$ & $1$   & $2$    & $5$    & $10$   & $20$\\\hline
$A(\rho)$& $1/6$ & $1/8$ & $1/12$ & $1/24$ & $1/44$ & $1/84$
\end{tabular}
\end{center}
Doubling $\rho$ from 1 to 2 requires 50\% more H-E compatibility density to
achieve the same loss level.

%----------------------------------------------------------------------
\section{Duality and Policy Comparisons}
\label{sec:duality}

\begin{lemma}[$\rho$-generalized Patient--TA-Greedy duality]
\label{lem:duality-rho}
\proved
\[
  L^{\mathrm{Patient}}(d_1,\rho)
  \;=\;
  L^{\mathrm{TA\text{-}Greedy}}\!\left(\frac{d_1}{2\rho}\right).
\]
\end{lemma}

\begin{proof}
From Lemma~\ref{lem:patient-rho}:
$L^{\mathrm{Patient}}(d_1,\rho)=4\rho W(d_1/(8\rho))/d_1$.
From Section~\ref{sec:TAG-rho}: for any $d_1'>0$,
$L^{\mathrm{TA\text{-}Greedy}}(d_1')=2W(d_1'/4)/d_1'$.
Setting $d_1'=d_1/(2\rho)$:
\[
  L^{\mathrm{TA\text{-}Greedy}}\!\left(\frac{d_1}{2\rho}\right)
  = \frac{2W\!\bigl(d_1/(8\rho)\bigr)}{d_1/(2\rho)}
  = \frac{4\rho\,W(d_1/(8\rho))}{d_1}
  = L^{\mathrm{Patient}}(d_1,\rho). \qedhere
\]
\end{proof}

\begin{remark}[Economic interpretation]
Outside-option pressure on E (higher $\rho$) is equivalent, for Patient
policy, to reducing the H-E compatibility density by a factor of $2\rho$.
A Patient planner facing E departure rate $\rho$ achieves the same loss as a
TA-Greedy planner in a market where $d_1$ has been divided by $2\rho$.
Conversely, halving $\rho$ delivers the same loss improvement as doubling
$d_1$.
\end{remark}

\paragraph{Policy ranking.}
For large $d_1$ (thick, heterogeneous markets), the loss ordering is:
\[
  L^{\text{TA-Patient}}
  \;\ll\;
  L^{\text{TA-Greedy}} \;\approx\; L^{\text{Greedy}}
  \;<\;
  L^{\text{Patient}},
\]
with TA-Patient achieving exponential loss and the others polynomial.
Importantly, \emph{without type-aware selection, being Patient is worse than
being Greedy}: under type-unaware Patient, the abundant E pool is consumed
by E-E matches at criticality rather than by H rescues.  Patience helps only
when the matching rule deliberately directs critical E agents toward H.

For thin markets ($d_1$ small), the ranking can differ.  In particular,
type-unaware Patient can be dominated by Greedy because waiting does not
build enough matching opportunities for H---the H pool is too sparse for
Patient's E-rescue mechanism to operate.

\paragraph{Threshold $\rho^*$.}
Numerical simulations (Figure: policy performance vs.\ E departure rate)
reveal a threshold $\rho^*(d_1,d_2)$ above which the Patient and TA-Patient
advantage over Greedy-type policies vanishes.  For $\rho>\rho^*$, E turnover
is fast enough that the E inventory is too thin to benefit from waiting.
Characterizing $\rho^*$ analytically is an open problem.

%----------------------------------------------------------------------
\section{Competition Between Exchanges}
\label{sec:competition}

\subsection{Two-platform model}

Consider two platforms serving the same population of H agents but
competing for E agents:
\begin{itemize}
\item \textbf{Platform A} (public): uses Patient policy; attracts fraction
  $\alpha\in(0,1)$ of E agents; E departure rate $\rho_0$.
\item \textbf{Platform B} (private): uses TA-Patient policy; attracts
  fraction $1-\alpha$ of E agents; E departure rate $\rho_B<\rho_0$
  (reduced by platform-specific retention mechanisms).
\end{itemize}
Platform B's ability to offer $\rho_B<\rho_0$ reflects tools unavailable to
the public platform: priority credits for altruistic donors, selective
offering of younger or better-matched recipients, and commitment mechanisms
that bind donors to the pool.

The composite E departure rate seen by H agents across both pools is
approximated by the harmonic mean:
\begin{equation}\label{eq:eff-rho}
  \rho_{\mathrm{eff}}(\alpha)
  = \left[\frac{\alpha}{\rho_0}+\frac{1-\alpha}{\rho_B}\right]^{-1}.
\end{equation}
As $\alpha$ decreases (more E agents migrate to the private platform),
$\rho_{\mathrm{eff}}$ falls: the system-wide E pool turns over more slowly.

\subsection{Welfare comparison}

\begin{proposition}[Competition can be welfare-improving]
\label{prop:competition}
\open For large $d_1$:
\[
  L^{\mathrm{TA\text{-}Patient}}\bigl(d_1,\,\rho_{\mathrm{eff}}(\alpha)\bigr)
  \;<\;
  L^{\mathrm{Patient}}(d_1,\rho_0),
\]
because exponential decay (TA-Patient) dominates polynomial decay (Patient)
as $d_1\to\infty$, even after the fragmentation penalty from $\alpha<1$.
For moderate $d_1$, welfare depends on a threshold
$d_1^*(\rho_0,\rho_B,\alpha)$.
\end{proposition}

\begin{remark}[Empirical context]
In the US kidney exchange system, the NKR (daily), the Alliance for Paired
Kidney Donation (daily or on new-entry), and the OPTN clearinghouse (weekly)
operate at different clearing frequencies, which correspond to different
effective E departure rates.  Hospital programs like UC's that match locally
before accessing national pools effectively run TA-Patient-like protocols.
Our model predicts that this layered, fragmented structure need not be
welfare-reducing: if faster-clearing platforms also implement type-aware
protocols and invest in donor retention (lower $\rho$), the TA-Patient
exponential advantage can outweigh the cost of pool fragmentation in thick
enough markets.
\end{remark}

%----------------------------------------------------------------------
\section{Discussion}
\label{sec:discussion}

\paragraph{Arrival imbalance.}
When E arrivals exceed H arrivals (fraction $\lambda_E>1/2$ of total
arrivals), the TA-Patient advantage is amplified: the deeper E pool
strengthens the rescue channel for H agents.  Numerical results show that
sufficient imbalance ($\lambda_E$ close to 1) can preserve the TA-Patient
advantage even at high $\rho$, because a thicker E pool compensates for
rapid E turnover.  This suggests a complementary design lever: platforms can
sustain the TA-Patient advantage under high $\rho$ by actively recruiting
disproportionately more E-type participants.

\paragraph{Finite-$d_2$ corrections.}
The $O(1/d_2)$ corrections to Greedy and TA-Greedy do depend on $\rho$:
when $d_2$ is finite, E-E matching becomes imperfect and E-pool dynamics
interact with $\rho$.  The crossover density $d_2^*$ at which policy
rankings flip in $\rho$ is an open numerical question.

\paragraph{Formal competition model.}
The two-platform welfare comparison relies on the effective-rate
approximation~\eqref{eq:eff-rho}, which treats the composite system as a
single pool with blended E departure rate.  A rigorous treatment would model
the joint Markov chain on $(H_t^A,E_t^A,H_t^B,E_t^B)$ and derive the coupled
mean-field fixed-point equations.

\paragraph{Stationary concentration.}
The mean-field approach requires the stationary distribution of the CTMC to
concentrate near the fixed point $(x^*,y^*)$.  This is supported numerically
by occupancy heatmaps that show tight mass near $(h^*,e^*)$, and is an
extension of the concentration arguments in the baseline AKL paper.

%----------------------------------------------------------------------
\section{Conclusion}
\label{sec:conclusion}

The E departure rate $\rho$ introduces a fundamental asymmetry across
policies: Patient-type policies are $\rho$-sensitive while Greedy-type
policies are not at leading order.  The closed form
$L^{\mathrm{Patient}}(d_1,\rho)=4\rho W(d_1/(8\rho))/d_1$ makes this
precise, and the duality $L^{\mathrm{Patient}}(d_1,\rho)=L^{\mathrm{TA\text{-}Greedy}}(d_1/(2\rho))$
translates it into a concrete equivalence: a unit increase in $\log\rho$ is
as damaging to Patient performance as a factor-$2\rho$ decrease in H-E
compatibility density.

For TA-Patient---the exponentially superior policy---higher $\rho$ compresses
the decay exponent from $d_1/8$ (at $\rho=1$) to $d_1/(4(1+\rho))$.  A
platform that retains E agents gains exponentially in thick markets.  The
competition model shows that fragmentation is not uniformly harmful: a
private platform that sufficiently reduces $\rho$ and uses TA-Patient matching
can dominate a public Patient platform once markets are thick enough.

\paragraph{Retention incentives.}
The analysis points to three concrete mechanisms for reducing $\rho$:
\begin{enumerate}
\item \emph{Priority credits}: E agents who donate accumulate credits for
  priority matching in the future, making staying in the exchange more
  attractive than outside options.
\item \emph{Quality matching}: Private exchanges can offer younger or
  healthier recipients to altruistic donors---improving the inside option
  relative to the outside option and directly reducing $\rho$.
\item \emph{Commitment mechanisms}: Soft commitments (listing the donor as
  active) or hard commitments (binding agreements with the exchange) reduce
  turnover directly.
\end{enumerate}
Each of these shifts the effective $\rho$ downward, and our results quantify
exactly how much: reducing $\rho$ by a factor of 2 is equivalent to doubling
$d_1$, i.e., halving the fraction of incompatible pairs in the market.

\end{document}
